% Options for packages loaded elsewhere
% Options for packages loaded elsewhere
\PassOptionsToPackage{unicode}{hyperref}
\PassOptionsToPackage{hyphens}{url}
\PassOptionsToPackage{dvipsnames,svgnames,x11names}{xcolor}
%
\documentclass[
  spanish,
  letterpaper,
  DIV=11,
  numbers=noendperiod]{scrartcl}
\usepackage{xcolor}
\usepackage{amsmath,amssymb}
\setcounter{secnumdepth}{-\maxdimen} % remove section numbering
\usepackage{iftex}
\ifPDFTeX
  \usepackage[T1]{fontenc}
  \usepackage[utf8]{inputenc}
  \usepackage{textcomp} % provide euro and other symbols
\else % if luatex or xetex
  \usepackage{unicode-math} % this also loads fontspec
  \defaultfontfeatures{Scale=MatchLowercase}
  \defaultfontfeatures[\rmfamily]{Ligatures=TeX,Scale=1}
\fi
\usepackage{lmodern}
\ifPDFTeX\else
  % xetex/luatex font selection
\fi
% Use upquote if available, for straight quotes in verbatim environments
\IfFileExists{upquote.sty}{\usepackage{upquote}}{}
\IfFileExists{microtype.sty}{% use microtype if available
  \usepackage[]{microtype}
  \UseMicrotypeSet[protrusion]{basicmath} % disable protrusion for tt fonts
}{}
\makeatletter
\@ifundefined{KOMAClassName}{% if non-KOMA class
  \IfFileExists{parskip.sty}{%
    \usepackage{parskip}
  }{% else
    \setlength{\parindent}{0pt}
    \setlength{\parskip}{6pt plus 2pt minus 1pt}}
}{% if KOMA class
  \KOMAoptions{parskip=half}}
\makeatother
% Make \paragraph and \subparagraph free-standing
\makeatletter
\ifx\paragraph\undefined\else
  \let\oldparagraph\paragraph
  \renewcommand{\paragraph}{
    \@ifstar
      \xxxParagraphStar
      \xxxParagraphNoStar
  }
  \newcommand{\xxxParagraphStar}[1]{\oldparagraph*{#1}\mbox{}}
  \newcommand{\xxxParagraphNoStar}[1]{\oldparagraph{#1}\mbox{}}
\fi
\ifx\subparagraph\undefined\else
  \let\oldsubparagraph\subparagraph
  \renewcommand{\subparagraph}{
    \@ifstar
      \xxxSubParagraphStar
      \xxxSubParagraphNoStar
  }
  \newcommand{\xxxSubParagraphStar}[1]{\oldsubparagraph*{#1}\mbox{}}
  \newcommand{\xxxSubParagraphNoStar}[1]{\oldsubparagraph{#1}\mbox{}}
\fi
\makeatother


\usepackage{longtable,booktabs,array}
\usepackage{calc} % for calculating minipage widths
% Correct order of tables after \paragraph or \subparagraph
\usepackage{etoolbox}
\makeatletter
\patchcmd\longtable{\par}{\if@noskipsec\mbox{}\fi\par}{}{}
\makeatother
% Allow footnotes in longtable head/foot
\IfFileExists{footnotehyper.sty}{\usepackage{footnotehyper}}{\usepackage{footnote}}
\makesavenoteenv{longtable}
\usepackage{graphicx}
\makeatletter
\newsavebox\pandoc@box
\newcommand*\pandocbounded[1]{% scales image to fit in text height/width
  \sbox\pandoc@box{#1}%
  \Gscale@div\@tempa{\textheight}{\dimexpr\ht\pandoc@box+\dp\pandoc@box\relax}%
  \Gscale@div\@tempb{\linewidth}{\wd\pandoc@box}%
  \ifdim\@tempb\p@<\@tempa\p@\let\@tempa\@tempb\fi% select the smaller of both
  \ifdim\@tempa\p@<\p@\scalebox{\@tempa}{\usebox\pandoc@box}%
  \else\usebox{\pandoc@box}%
  \fi%
}
% Set default figure placement to htbp
\def\fps@figure{htbp}
\makeatother



\ifLuaTeX
\usepackage[bidi=basic]{babel}
\else
\usepackage[bidi=default]{babel}
\fi
% get rid of language-specific shorthands (see #6817):
\let\LanguageShortHands\languageshorthands
\def\languageshorthands#1{}


\setlength{\emergencystretch}{3em} % prevent overfull lines

\providecommand{\tightlist}{%
  \setlength{\itemsep}{0pt}\setlength{\parskip}{0pt}}



 


\KOMAoption{captions}{tableheading}
\makeatletter
\@ifpackageloaded{caption}{}{\usepackage{caption}}
\AtBeginDocument{%
\ifdefined\contentsname
  \renewcommand*\contentsname{Tabla de contenidos}
\else
  \newcommand\contentsname{Tabla de contenidos}
\fi
\ifdefined\listfigurename
  \renewcommand*\listfigurename{Listado de Figuras}
\else
  \newcommand\listfigurename{Listado de Figuras}
\fi
\ifdefined\listtablename
  \renewcommand*\listtablename{Listado de Tablas}
\else
  \newcommand\listtablename{Listado de Tablas}
\fi
\ifdefined\figurename
  \renewcommand*\figurename{Figura}
\else
  \newcommand\figurename{Figura}
\fi
\ifdefined\tablename
  \renewcommand*\tablename{Tabla}
\else
  \newcommand\tablename{Tabla}
\fi
}
\@ifpackageloaded{float}{}{\usepackage{float}}
\floatstyle{ruled}
\@ifundefined{c@chapter}{\newfloat{codelisting}{h}{lop}}{\newfloat{codelisting}{h}{lop}[chapter]}
\floatname{codelisting}{Listado}
\newcommand*\listoflistings{\listof{codelisting}{Listado de Listados}}
\makeatother
\makeatletter
\makeatother
\makeatletter
\@ifpackageloaded{caption}{}{\usepackage{caption}}
\@ifpackageloaded{subcaption}{}{\usepackage{subcaption}}
\makeatother
\usepackage{bookmark}
\IfFileExists{xurl.sty}{\usepackage{xurl}}{} % add URL line breaks if available
\urlstyle{same}
\hypersetup{
  pdftitle={Fundamentación Teórica y Formulación Matemática del Modelo},
  pdflang={es},
  colorlinks=true,
  linkcolor={blue},
  filecolor={Maroon},
  citecolor={Blue},
  urlcolor={Blue},
  pdfcreator={LaTeX via pandoc}}


\title{Fundamentación Teórica y Formulación Matemática del Modelo}
\author{}
\date{}
\begin{document}
\maketitle


\subsection{Introducción al problema integrado de localización e
inventario}\label{introducciuxf3n-al-problema-integrado-de-localizaciuxf3n-e-inventario}

\subsubsection{Relevancia en logística humanitaria y gestión de
emergencias}\label{relevancia-en-loguxedstica-humanitaria-y-gestiuxf3n-de-emergencias}

La logística humanitaria constituye un componente esencial para la
respuesta a desastres, pues se orienta a garantizar la disponibilidad,
distribución y acceso equitativo a bienes vitales en contextos
caracterizados por alta incertidumbre, infraestructura deteriorada y
restricciones severas de recursos (Balcik \& Beamon, 2008; Caunhye et
al., 2012). A diferencia de los sistemas logísticos comerciales, donde
predominan objetivos de eficiencia, los entornos humanitarios priorizan
rapidez, cobertura y reducción del riesgo, lo que exige modelos
matemáticos capaces de integrar simultáneamente decisiones estratégicas
y tácticas bajo incertidumbre.

En este marco, el \textbf{problema integrado de localización e
inventario} se ha consolidado como una herramienta fundamental para la
planificación anticipada, ya que permite determinar, de manera conjunta:

\begin{enumerate}
\def\labelenumi{\arabic{enumi}.}
\tightlist
\item
  \textbf{La ubicación óptima de almacenes de preposicionamiento}, y\\
\item
  \textbf{Las cantidades de inventario que deben almacenarse antes del
  evento disruptivo}.
\end{enumerate}

Ambas decisiones condicionan el desempeño global de la red en términos
de costo total esperado, robustez ante escenarios adversos y equidad en
la atención (Vanajakum \& Kachitvichyanukul, 2023). Su naturaleza dual
---estratégica y táctica simultáneamente--- explica su complejidad
estructural.

\subsubsection{Formulación matemática
general}\label{formulaciuxf3n-matemuxe1tica-general}

Sean los conjuntos:

\begin{itemize}
\tightlist
\item
  \(\mathcal{I} = \{1,\dots,m\}\): ubicaciones candidatas para
  almacenes,\\
\item
  \(\mathcal{J} = \{1,\dots,n\}\): zonas de demanda,\\
\item
  \(d_j \sim \mathcal{D}_j\): demanda aleatoria en la zona \(j\),\\
\item
  \(x_i \in \{0,1\}\): variable de apertura en la ubicación \(i\),\\
\item
  \(Q_i \ge 0\): inventario preposicionado,\\
\item
  \(c_{ij}\): costo unitario de transporte,\\
\item
  \(h_i\): costo unitario de mantenimiento de inventario,\\
\item
  \(\beta_j\): penalización por faltante en la zona \(j\).
\end{itemize}

El modelo se estructura como un \textbf{programa estocástico de dos
etapas}, donde la primera etapa decide \(x\) y \(Q\), mientras que la
segunda etapa determina los envíos \(y_{ij}(\xi)\) según la realización
aleatoria \(\xi\):

\[
\min_{x \in \{0,1\}^m,\; Q \ge 0}
\left[
\sum_{i \in \mathcal{I}} f_i(x_i)
+ \sum_{i \in \mathcal{I}} h_i Q_i
+ \mathbb{E}_\xi \left(
\min_{y(\xi)\ge 0}
\left\{
\sum_{i \in \mathcal{I}} \sum_{j \in \mathcal{J}} c_{ij}\, y_{ij}(\xi)
+ \sum_{j \in \mathcal{J}} \beta_j
\left( d_j(\xi) - \sum_{i \in \mathcal{I}} y_{ij}(\xi) \right)^+
\right\}
\right)
\right].
\]

Sujeto a las restricciones de demanda y capacidad:

\[
\sum_{i \in \mathcal{I}} y_{ij}(\xi) \le d_j(\xi), \qquad \forall j, \; \forall \xi,
\]

\[
\sum_{j \in \mathcal{J}} y_{ij}(\xi) \le Q_i, \qquad \forall i, \; \forall \xi,
\]

\[
Q_i \le M x_i, \qquad \forall i,
\]

donde \((\cdot)^+ = \max\{0,\cdot\}\) y \(M\) es un parámetro de
activación suficientemente grande.

Este modelo pertenece a la clase \textbf{SMINLP} (Stochastic
Mixed-Integer Nonlinear Programs), cuya solución requiere técnicas
computacionales avanzadas como descomposición de Benders estocástica,
aproximación por escenarios o muestreo tipo Monte Carlo (Shapiro et al.,
2021).

\subsubsection{Ejemplo ilustrativo (caso determinista
simplificado)}\label{ejemplo-ilustrativo-caso-determinista-simplificado}

Para ilustrar la interacción entre localización, inventario y
asignación, considérese:

\begin{itemize}
\tightlist
\item
  \(m = 2\) almacenes,\\
\item
  \(n = 2\) zonas de demanda,\\
\item
  demandas deterministas: \(d_1 = d_2 = 40\),\\
\item
  costos fijos: \(f_1 = 500\), \(f_2 = 600\),\\
\item
  costos de transporte:

  \begin{itemize}
  \tightlist
  \item
    \(c_{11} = 10\), \(c_{12} = 20\),\\
  \item
    \(c_{21} = 15\), \(c_{22} = 12\),\\
  \end{itemize}
\item
  costos de mantenimiento: \(h_1 = h_2 = 2\),\\
\item
  penalización por faltante: \(\beta_1 = \beta_2 = 100\).
\end{itemize}

\paragraph{Caso 1. Abrir solo el almacén
1}\label{caso-1.-abrir-solo-el-almacuxe9n-1}

Inventario mínimo requerido: \(Q_1 = 80\).

Costo total:

\[
500 + 2(80) + 10(40) + 20(40)
= 500 + 160 + 400 + 800 = 1860.
\]

\paragraph{Caso 2. Abrir solo el almacén
2}\label{caso-2.-abrir-solo-el-almacuxe9n-2}

Costo total:

\[
600 + 2(80) + 15(40) + 12(40)
= 600 + 160 + 600 + 480 = 1840.
\]

\paragraph{Caso 3. Abrir ambos
almacenes}\label{caso-3.-abrir-ambos-almacenes}

Asignación eficiente:

\begin{itemize}
\tightlist
\item
  Almacén 1 → Zona 1: 40 unidades,\\
\item
  Almacén 2 → Zona 2: 40 unidades,\\
\item
  \(Q_1 = Q_2 = 40\).
\end{itemize}

Costo total:

\[
500 + 600 + 2(40) + 2(40) + 400 + 480
= 1100 + 160 + 880 = 2140.
\]

La mejor decisión es abrir únicamente el almacén 2.\\
Este resultado evidencia que:

\begin{itemize}
\tightlist
\item
  abrir más almacenes no siempre reduce el costo total,\\
\item
  la estructura espacial del costo de transporte domina la decisión,\\
\item
  en presencia de incertidumbre, la solución óptima puede favorecer
  redundancia para reducir riesgo.
\end{itemize}

Lo anterior justifica la necesidad de modelos estocásticos e integrados
en la planificación humanitaria contemporánea.

\subsubsection{Limitaciones de los enfoques clásicos bajo incertidumbre
y no
linealidad}\label{limitaciones-de-los-enfoques-cluxe1sicos-bajo-incertidumbre-y-no-linealidad}

Los modelos clásicos de localización e inventario ---como el problema de
localización sin capacidades (UFLP), el modelo de lote económico (EOQ) o
sus variantes deterministas--- han sido ampliamente estudiados bajo el
supuesto de \textbf{información perfecta y estructura lineal}, un
enfoque cuya pertinencia disminuye notablemente en escenarios
humanitarios caracterizados por \textbf{incertidumbre estructural} y
\textbf{no linealidad inherente} en costos y restricciones (Hadley \&
Whitin, 1963; Snyder, 2006).

\paragraph{Supuestos restrictivos del modelo EOQ
clásico}\label{supuestos-restrictivos-del-modelo-eoq-cluxe1sico}

El modelo EOQ se fundamenta en supuestos como:

\begin{enumerate}
\def\labelenumi{\arabic{enumi}.}
\tightlist
\item
  Demanda constante y conocida \(D>0\).
\item
  Tiempo de entrega determinista \(L \geq 0\).
\item
  Costos lineales: costo de pedido \(S\), costo unitario de
  mantenimiento \(H\), sin faltantes.
\item
  Reposición instantánea y sin restricciones de capacidad (Hadley \&
  Whitin, 1963).
\end{enumerate}

El costo total del modelo es:

\[
Z^{EOQ}(Q)=\frac{D}{Q}S+\frac{Q}{2}H+DC,
\]

y su solución óptima clásica:

\[
Q^{EOQ}=\sqrt{\frac{2DS}{H}}.
\]

Este resultado depende de la \textbf{deterministicidad} y la
\textbf{convexidad estricta}. Sin embargo, en logística humanitaria:

\begin{itemize}
\tightlist
\item
  La demanda es aleatoria \(D(\omega)\).
\item
  \(L\) es incierto debido a afectaciones en infraestructura.
\item
  Los faltantes deben modelarse mediante una penalización finita
  \(\beta>0\).
\item
  Los costos pueden presentar economías de escala y no convexidades
  (Porteus, 2002).
\end{itemize}

\paragraph{Incapacidad del EOQ para capturar
riesgo}\label{incapacidad-del-eoq-para-capturar-riesgo}

Si la demanda durante el tiempo de entrega es
\(D_L\sim\mathcal{N}(\mu_L,\sigma_L^2)\), el inventario de seguridad
requiere:

\[
\mathbb{P}(D_L\leq Q/2 + Z_\alpha\sigma_L)\geq \alpha.
\]

Con ello, la función de costo extendida se vuelve:

\[
Z(Q)=\frac{D}{Q}S+\left(\frac{Q}{2}+Z_\alpha\sigma_L\right)H+DC+\frac{D}{Q}\beta C. \tag{2.1}
\]

Esta función \textbf{no es convexa globalmente} cuando \(\sigma_L\)
depende de \(Q\), como ocurre cuando existe congestión en transporte o
variabilidad dependiente del tamaño del pedido (Bertsimas \& Thiele,
2006; Shapiro et al., 2021).

\begin{quote}
\textbf{Proposición 2.1.2.1 (No linealidad estricta).}\\
La función \(Z(Q)\) definida en (2.1) es diferenciable en
\((0,\infty)\), pero no es convexa globalmente si \(\beta>0\) y
\(\sigma_L>0\).\\
Si \(\sigma_L\) es constante: \[
\frac{d^2 Z}{dQ^2}=\frac{2D(S+\beta C)}{Q^3}>0.
\] Pero si \(\sigma_L=\sigma\sqrt{L(Q)}\) y \(L(Q)\) es no lineal,
entonces \(Z(Q)\) pierde convexidad.
\end{quote}

Esto demuestra que muchos métodos clásicos basados en convexidad y
separabilidad resultan inaplicables en condiciones realistas.

\paragraph{Limitaciones de modelos de localización
deterministas}\label{limitaciones-de-modelos-de-localizaciuxf3n-deterministas}

El modelo clásico UFLP resuelve:

\[
\min_{x\in\{0,1\}^m, y\ge 0}\sum_{i}f_i x_i + \sum_{i}\sum_j c_{ij}y_{ij},
\]

bajo restricciones de asignación deterministas. Sin embargo:

\begin{itemize}
\tightlist
\item
  La demanda \(d_j\) es aleatoria.
\item
  Los costos pueden variar según accesibilidad o daños.
\item
  Las capacidades reales son inciertas.
\end{itemize}

Modelar \(d_j\) mediante \(\mathbb{E}[d_j]\) genera soluciones frágiles
frente a escenarios adversos (Snyder, 2006).

\paragraph{Ejemplo ilustrativo: EOQ vs.~modelo
estocástico}\label{ejemplo-ilustrativo-eoq-vs.-modelo-estocuxe1stico}

Considérese:

\begin{itemize}
\tightlist
\item
  \(D=1000\) pero en realidad \(D\sim\text{Uniform}(800,1200)\),
\item
  \(S=100\), \(H=5\), \(C=10\), \(\beta=50\),
\item
  \(L=0.1\), \(\sigma=10 \Rightarrow \sigma_L\approx 3.16\),
\item
  Nivel de servicio \(95\% \Rightarrow Z_\alpha=1.645\).
\end{itemize}

\textbf{EOQ clásico:}

\[
Q^{EOQ}=\sqrt{\frac{2(1000)(100)}{5}}=200,
\]

y

\[
Z^{EOQ}=11000.
\]

\textbf{Modelo estocástico:}

\[
Z(Q)=\frac{600000}{Q}+2.5Q+10025.99.
\]

Óptimo:

\[
Q^*=\sqrt{\frac{600000}{2.5}}\approx 489.9,
\]

y

\[
Z(Q^*)\approx 12475.44.
\]

Si se usa \(Q=200\):

\[
Z(200)=13525.99,
\]

lo que implica un \textbf{8.4\% más de costo}, además de riesgos de
faltante significativamente mayores.

Este contraste demuestra que los modelos deterministas pueden producir
resultados \textbf{subóptimos o inadmisibles} en contextos humanitarios
donde la demanda insatisfecha tiene consecuencias operativas y éticas
relevantes.

\paragraph{Conclusión de la
subsección}\label{conclusiuxf3n-de-la-subsecciuxf3n}

Los enfoques clásicos, por su naturaleza determinista, \textbf{no
incorporan mecanismos para manejar riesgo, variabilidad ni no
linealidades estructurales}. Por ello, en logística humanitaria se
requiere un marco que:

\begin{itemize}
\tightlist
\item
  integre incertidumbre probabilística;
\item
  incorpore penalizaciones no lineales;
\item
  respete la naturaleza discreta-continua del problema;
\item
  mantenga propiedades analíticas bajo condiciones realistas.
\end{itemize}

Este es precisamente el propósito del modelo desarrollado en este
trabajo.

\subsubsection{Estructura híbrida: decisiones discretas (ubicación) y
continuas
(inventario)}\label{estructura-huxedbrida-decisiones-discretas-ubicaciuxf3n-y-continuas-inventario}

Los problemas integrados de localización e inventario forman parte de
una clase fundamental denominada \emph{modelos híbridos} o
\emph{problemas de decisión mixta}, en los que coexisten variables
discretas y continuas. Este tipo de modelos presenta una complejidad
inherente debida a la ruptura topológica inducida por las variables
binarias, lo cual conlleva dominios no convexos y no conexos, impidiendo
la aplicación directa de técnicas clásicas basadas en convexidad
(Wolsey, 1998; Grossmann, 2002; Daskin, 2013).

\paragraph{Formulación matemática del espacio de
decisiones}\label{formulaciuxf3n-matemuxe1tica-del-espacio-de-decisiones}

Sea: - \(\mathcal{I} = \{1,\ldots,m\}\) el conjunto de ubicaciones
posibles para almacenes, - \(\mathcal{J} = \{1,\ldots,n\}\) las zonas de
demanda, - \(x = (x_1,\dots,x_m)^\top \in \{0,1\}^m\) las decisiones
binarias de apertura, - \(y = (y_{ij}) \in \mathbb{R}_+^{m\times n}\)
los flujos desde almacenes a zonas, -
\(s = (s_1,\dots,s_m)^\top \in \mathbb{R}_+^m\) el inventario
preposicionado.

El espacio factible se define como:

\[
\mathcal{X} := 
\left\{
(x,y,s)\in \{0,1\}^m\times \mathbb{R}_+^{m\times n}\times \mathbb{R}_+^m \;\middle|\;
\begin{aligned}
&\sum_{i\in\mathcal{I}} y_{ij} \ge d_j, &&\forall j\in\mathcal{J},\\
&\sum_{j\in\mathcal{J}} y_{ij} \le s_i, &&\forall i\in\mathcal{I},\\
&s_i \le M x_i, &&\forall i\in\mathcal{I}
\end{aligned}
\right\}.
\tag{2.2}
\]

Aquí, \(d_j > 0\) representa la demanda de la zona \(j\) y \(M\) es un
parámetro de gran magnitud. La tercera restricción describe un vínculo
lógico: si \(x_i = 0\), entonces \(s_i = 0\) y, en consecuencia,
\(y_{ij} = 0\). Este acoplamiento induce una \textbf{no convexidad
estructural} (Bertsekas, 1999).

\begin{quote}
\textbf{Proposición (No convexidad).}\\
El conjunto \(\mathcal{X}\) definido en (2.2) \textbf{no es convexo}.\\
\emph{Demostración.} Considérense las configuraciones factibles
\(x^{(1)}=(1,0,\dots,0)\) y \(x^{(2)}=(0,1,0,\dots,0)\).\\
El punto medio \(\tfrac12 x^{(1)} + \tfrac12 x^{(2)}\) tiene componentes
iguales a \(1/2\), lo cual viola la condición binaria. Luego,
\(\mathcal{X}\) no es convexo. ∎
\end{quote}

La consecuencia directa es que condiciones de optimalidad de tipo KKT
\textbf{no garantizan} optimalidad global (Lewis \& Overton, 2013).

\begin{center}\rule{0.5\linewidth}{0.5pt}\end{center}

\paragraph{Función objetivo y acoplamiento
estructural}\label{funciuxf3n-objetivo-y-acoplamiento-estructural}

El costo total del sistema se modela como:

\[
Z(x,y,s) 
= 
\underbrace{\sum_{i\in\mathcal{I}} f_i(x_i)}_{\text{Costos fijos}}
+ \underbrace{\sum_{i\in\mathcal{I}} h_i s_i}_{\text{Costos de inventario}}
+ \underbrace{\sum_{i\in\mathcal{I}}\sum_{j\in\mathcal{J}} c_{ij} y_{ij}}_{\text{Costos de transporte}}.
\tag{2.3}
\]

Aunque la parte en variables continuas es lineal, el término
\(f_i(x_i)\) genera una \textbf{no linealidad discreta}; más aún, en la
práctica humanitaria pueden existir economías de escala, accesibilidad
variable o costos de apertura no lineales, reforzando la complejidad
global (Grossmann, 2002).

El sistema presenta una jerarquía:

\begin{enumerate}
\def\labelenumi{\arabic{enumi}.}
\tightlist
\item
  \(x\) determina la infraestructura disponible,\\
\item
  \(s\) depende de \(x\),\\
\item
  \(y\) depende de \(s\).
\end{enumerate}

Esto produce una dependencia en cascada que elimina la separabilidad
global del problema.

\begin{quote}
\textbf{Proposición (Convexidad condicional).}\\
Para cualquier \(\bar{x} \in \{0,1\}^m\), el subproblema

\[
\min_{(y,s) \in \mathcal{X}(\bar{x})} Z(\bar{x}, y, s)
\]

es un programa lineal y, por tanto, convexo.\\
\emph{Demostración.} Fijado \(\bar{x}\), todas las restricciones son
lineales en \((y,s)\) y la función objetivo es afín. ∎
\end{quote}

Esta propiedad respalda el uso de descomposición tipo Benders o esquemas
maestro--subproblema (Wolsey, 1998).

\begin{center}\rule{0.5\linewidth}{0.5pt}\end{center}

\paragraph{Ejemplo ilustrativo: interacción
discreto--continuo}\label{ejemplo-ilustrativo-interacciuxf3n-discretocontinuo}

Considérese: - \(m=2\) almacenes,\\
- \(n=2\) demandas con \(d_1=d_2=50\),\\
- Costos fijos \(f_1=300\), \(f_2=400\),\\
- Costos de inventario \(h_1=h_2=3\),\\
- Costos de transporte:\\
\(c_{11}=8\), \(c_{12}=15\),\\
\(c_{21}=12\), \(c_{22}=7\),\\
- \(M=100\).

\textbf{(i) Abrir solo 1:}\\
\[
Z = 300 + 3\cdot 100 + (8\cdot 50 + 15\cdot 50)=1750.
\]

\textbf{(ii) Abrir solo 2:}\\
\[
Z = 400 + 3\cdot 100 + (12\cdot 50 + 7\cdot 50)=1650.
\]

\textbf{(iii) Abrir ambos:}\\
Asignación eficiente, pero mayor costo fijo:\\
\[
Z = 300 + 400 + 3\cdot 100 + (8\cdot 50 + 7\cdot 50)=1750.
\]

Conclusión: es óptimo \textbf{abrir solo el almacén 2}.\\
La solución cambiaría si los costos fijos disminuyeran, lo que evidencia
la \emph{sensibilidad estructural} del diseño híbrido (Daskin, 2013).

La integración simultánea de decisiones discretas y continuas induce
dominios no convexos, dependencias lógicas y estructuras no separables
que requieren marcos teóricos propios de la programación entera mixta,
análisis no suave y descomposición matemática. Esta base establece los
fundamentos del modelo propuesto en las siguientes secciones, donde se
formaliza la integración bajo condiciones logísticas realistas y
restricciones humanitarias.

\subsubsection{\texorpdfstring{Conjuntos de índices: instalaciones
potenciales \$ \mathcal{I} \$ y zonas de demanda \$ \mathcal{J}
\$}{Conjuntos de índices: instalaciones potenciales \$  \$ y zonas de demanda \$  \$}}\label{conjuntos-de-uxedndices-instalaciones-potenciales-y-zonas-de-demanda}

La construcción rigurosa de un modelo de optimización comienza con la
definición precisa de los conjuntos de índices, los cuales encapsulan la
estructura topológica y semántica del sistema bajo estudio. En el
contexto del problema integrado de localización e inventario para
logística humanitaria, dos conjuntos indexados fundamentales determinan
la arquitectura del modelo: el conjunto de instalaciones potenciales y
el conjunto de zonas de demanda.

\paragraph{\texorpdfstring{\textbf{Definición (Conjunto de instalaciones
potenciales).}}{Definición (Conjunto de instalaciones potenciales).}}\label{definiciuxf3n-conjunto-de-instalaciones-potenciales.}

Sea

\[
\mathcal{I} := \{1, 2, \dots, m\}, \quad m \in \mathbb{N},\; m \geq 1,
\]

el conjunto finito y numerable de ubicaciones candidatas para la
apertura de almacenes de preposicionamiento. Cada índice
\(i \in \mathcal{I}\) representa una localización geográfica factible
donde podría instalarse una instalación de almacenamiento y
distribución.

\paragraph{\texorpdfstring{\textbf{Definición (Conjunto de zonas de
demanda).}}{Definición (Conjunto de zonas de demanda).}}\label{definiciuxf3n-conjunto-de-zonas-de-demanda.}

Sea

\[
\mathcal{J} := \{1, 2, \dots, n\}, \quad n \in \mathbb{N},\; n \geq 1,
\]

el conjunto finito y numerable de zonas afectadas o potenciales de
demanda. Cada índice \(j \in \mathcal{J}\) denota una región geográfica
donde se anticipa la necesidad de suministros humanitarios tras un
evento disruptivo.

\begin{quote}
\textbf{Observación}\\
Los conjuntos \$ \mathcal{I} \$ y \$ \mathcal{J} \$ se asumen disjuntos
en función, aunque no necesariamente en localización física. Desde el
punto de vista del modelo, los roles son distintos:\\
- Los elementos de \$ \mathcal{I} \$ son fuentes de suministro.\\
- Los elementos de \$ \mathcal{J} \$ son sumideros de demanda.
\end{quote}

La cardinalidad de estos conjuntos ---es decir, \(m = |\mathcal{I}|\) y
\(n = |\mathcal{J}|\)--- determina directamente la complejidad
combinatoria del problema. En particular, el número de configuraciones
posibles de apertura de almacenes es \(2^m\), lo que implica que incluso
para valores modestos (por ejemplo, \(m = 30\)), el espacio discreto es
del orden de \(10^9\). Esta explosión combinatoria subraya la necesidad
de métodos de optimización estructurada.

\paragraph{\texorpdfstring{\textbf{Ejemplo
ilustrativo}}{Ejemplo ilustrativo}}\label{ejemplo-ilustrativo}

Supóngase un escenario de planificación para respuesta a inundaciones
dividido en cinco distritos. Se identifican tres ubicaciones potenciales
para almacenes:

\begin{itemize}
\tightlist
\item
  \(i = 1\): Aeropuerto internacional,\\
\item
  \(i = 2\): Base militar,\\
\item
  \(i = 3\): Centro de distribución de una ONG.
\end{itemize}

Asimismo, se consideran cuatro zonas vulnerables:

\begin{itemize}
\tightlist
\item
  \(j = 1\): Zona ribereña,\\
\item
  \(j = 2\): Área rural,\\
\item
  \(j = 3\): Barrio urbano denso,\\
\item
  \(j = 4\): Comunidad indígena aislada.
\end{itemize}

Los conjuntos índice resultan:

\[
\mathcal{I} = \{1,2,3\}, \quad 
\mathcal{J} = \{1,2,3,4\}.
\]

Sobre estos conjuntos se definirán:

\begin{itemize}
\tightlist
\item
  variables binarias \(x_i \in \{0,1\}\) para cada
  \(i \in \mathcal{I}\),\\
\item
  variables continuas de flujo \(y_{ij} \ge 0\) para
  \((i,j) \in \mathcal{I} \times \mathcal{J}\),\\
\item
  parámetros exógenos como \(c_{ij}\) (costos), \(f_i\) (costos fijos) y
  \(d_j\) (demanda).
\end{itemize}

\begin{quote}
\textbf{Nota metodológica.}\\
En modelado riguroso, los conjuntos índice deben definirse antes que las
variables, dado que estas heredan su dominio directamente de aquellos.
\end{quote}

Los conjuntos \$ \mathcal{I} \$ y \$ \mathcal{J} \$ constituyen
estructuras matemáticas esenciales que determinan la escala, la
conectividad y la complejidad del modelo. Su correcta definición resulta
crucial para la formalización posterior de variables, restricciones y
dependencias del sistema.

\subsubsection{\texorpdfstring{Variables de decisión: binarias (\(x_i\))
y continuas
(\(y_{ij}\))}{Variables de decisión: binarias (x\_i) y continuas (y\_\{ij\})}}\label{variables-de-decisiuxf3n-binarias-x_i-y-continuas-y_ij}

Una vez establecida la estructura combinatoria del problema mediante los
conjuntos índice \(\mathcal{I}\) y \(\mathcal{J}\), el siguiente paso en
la formulación rigurosa de un modelo de optimización es la definición
precisa del \textbf{espacio de decisiones}, mediante la especificación
de las \textbf{variables de decisión}. En el problema integrado de
localización e inventario, este espacio es híbrido, compuesto por dos
tipos cualitativamente distintos de variables:

\begin{itemize}
\tightlist
\item
  Variables binarias, que modelan elecciones discretas de
  infraestructura,
\item
  Variables continuas, que representan flujos físicos de recursos.
\end{itemize}

Esta dualidad refleja la naturaleza multinivel del problema: decisiones
estratégicas (¿dónde construir?) y decisiones tácticas u operativas
(¿cuánto enviar?).

\paragraph{\texorpdfstring{\textbf{Definición (Variables binarias de
localización).}}{Definición (Variables binarias de localización).}}\label{definiciuxf3n-variables-binarias-de-localizaciuxf3n.}

Para cada \(i \in \mathcal{I}\), se define la variable binaria

\[
x_i \in \{0,1\},
\]

donde

\[
x_i =
\begin{cases}
1, & \text{si se decide abrir una instalación en la ubicación } i, \\
0, & \text{en caso contrario}.
\end{cases}
\]

El vector \(x = (x_i)_{i \in \mathcal{I}} \in \{0,1\}^m\) codifica una
\textbf{configuración de infraestructura} del sistema logístico. Estas
variables pertenecen a un espacio discreto no convexo y finito cuya
cardinalidad es \(2^m\).

\paragraph{\texorpdfstring{\textbf{Definición (Variables continuas de
flujo).}}{Definición (Variables continuas de flujo).}}\label{definiciuxf3n-variables-continuas-de-flujo.}

Para cada par \((i,j) \in \mathcal{I} \times \mathcal{J}\), se define la
variable continua

\[
y_{ij} \in \mathbb{R}_+,
\]

que representa la cantidad de unidades asignadas desde la instalación
\(i\) a la zona de demanda \(j\) en el período de planificación.

El tensor
\(y = (y_{ij})_{(i,j) \in \mathcal{I} \times \mathcal{J}} \in \mathbb{R}_+^{m \times n}\)
describe la estructura de asignación del sistema, y pertenece a un
espacio vectorial convexo, cerrado y de dimensión finita.

\begin{quote}
\textbf{Observación (Naturaleza híbrida del espacio de decisiones).}\\
El espacio total de decisiones es el producto cartesiano\\
\[
\mathcal{D} = \{0,1\}^m \times \mathbb{R}_+^{m \times n},
\] que es no convexo, no conexo y discontinuo en la dirección discreta.
Esta estructura impide la aplicación directa de herramientas del
análisis convexo clásico y exige métodos de optimización mixta.
\end{quote}

Las variables \(x_i\) y \(y_{ij}\) no son independientes. Su
acoplamiento se establece mediante \textbf{restricciones de
factibilidad}, que garantizan que no se asignen flujos desde
instalaciones no abiertas. Para todo \(i \in \mathcal{I}\) y
\(j \in \mathcal{J}\):

\[
y_{ij} \le M x_i, \tag{2.4}
\]

donde \(M > 0\) es una constante suficientemente grande (por ejemplo,
\(M = \sum_{j \in \mathcal{J}} d_j\)). La restricción (2.4) es una
formulación \emph{big-M} estándar que modela la implicación lógica:

\[
x_i = 0 \;\Rightarrow\; y_{ij} = 0 \quad \forall j \in \mathcal{J}.
\]

\begin{center}\rule{0.5\linewidth}{0.5pt}\end{center}

\subsubsection{\texorpdfstring{\textbf{Ejemplo ilustrativo:
interpretación y acoplamiento de
variables}}{Ejemplo ilustrativo: interpretación y acoplamiento de variables}}\label{ejemplo-ilustrativo-interpretaciuxf3n-y-acoplamiento-de-variables}

Sea:

\[
\mathcal{I} = \{1,2,3\}, \qquad \mathcal{J} = \{1,2,3,4\}.
\]

Supóngase que se decide abrir instalaciones en \(i=1\) y \(i=3\), pero
no en \(i=2\):

\[
x = (1,0,1).
\]

Las restricciones big-M implican:

\begin{itemize}
\tightlist
\item
  \(y_{2j} = 0\) para todo \(j\),
\item
  \(y_{1j}, y_{3j} \ge 0\) como flujos posibles.
\end{itemize}

Una asignación factible es:

\[
\begin{aligned}
y_{11} &= 30, & y_{12} &= 20, & y_{13} &= 0, & y_{14} &= 0, \\
y_{21} &= y_{22} = y_{23} = y_{24} = 0, \\
y_{31} &= 10, & y_{32} &= 0, & y_{33} &= 40, & y_{34} &= 25.
\end{aligned}
\]

Si la demanda es \(d = (40,20,40,25)\), esta asignación es factible y
completa.

Este ejemplo muestra que las variables binarias \textbf{activan o
desactivan} subespacios del espacio continuo de flujos, generando una
partición de \(\mathcal{D}\) en \(2^m\) subespacios convexos.

Las variables \(x_i\) y \(y_{ij}\) constituyen los grados de libertad
fundamentales del modelo. Su naturaleza híbrida ---discreta y
continua--- refleja la dualidad entre planeación estratégica y operación
táctica. El acoplamiento lógico mediante restricciones big-M introduce
no convexidad, justificando el uso de técnicas de optimización entera y
métodos especializados. Las siguientes secciones integrarán estas
variables en la función objetivo y en las restricciones de balance de
flujo.

\subsubsection{Parámetros del sistema: costos fijos, unitarios, demanda
y
capacidades}\label{paruxe1metros-del-sistema-costos-fijos-unitarios-demanda-y-capacidades}

En la formulación matemática de un problema de optimización, los
\textbf{parámetros} constituyen los datos exógenos que estructuran el
espacio factible y la función objetivo. Su correcta especificación es
requisito previo a la definición de variables y restricciones; por
tanto, deben presentarse con dominio y unidades explícitas (práctica
recomendada por Lewis y colaboradores; Lewis, 2003). En el problema
integrado de localización e inventario consideramos, de modo explícito,
las cuatro familias de parámetros siguientes.

\paragraph{Costos fijos de apertura}\label{costos-fijos-de-apertura}

Para cada \(i\in\mathcal{I}\) definimos \[
f_i \in \mathbb{R}_{+}
\] como el \textbf{costo fijo de apertura} de la instalación en la
localización \(i\). En la formulación estándar el término contribuye a
la función objetivo mediante \(f_i x_i\) si \(x_i\) es la variable
binaria de apertura. Se exige \[
f_i \ge 0,\qquad \forall i\in\mathcal{I}.
\]

Observación práctica: cuando la apertura implica economías o
penalizaciones no triviales (por ejemplo, coste creciente con la
capacidad instalada o costes discretos adicionales por remoción),
\(f_i\) puede modelarse como función \(f_i(x_i,s_i)\); sin embargo, para
la formulación base se adopta linealidad en \(x_i\) y se externalizan
las no linealidades en extensiones (Shapiro et al., 2021).

\paragraph{Costos unitarios de
transporte}\label{costos-unitarios-de-transporte}

Para cada par \((i,j)\in\mathcal{I}\times\mathcal{J}\) definimos \[
c_{ij}\in\mathbb{R}_{+},
\] el \textbf{costo unitario de transportar} una unidad desde la
instalación \(i\) hasta la zona de demanda \(j\). Requerimos \[
c_{ij}\ge 0,\quad \forall i,j,
\] y denotamos la matriz
\(C:=[c_{ij}]_{i,j}\in\mathbb{R}_+^{m\times n}\). En contextos
humanitarios \(c_{ij}\) incorpora distancia, tiempo de acceso y
penalizaciones por rutas inseguras; su estimación suele derivar de
información geoespacial y modelos de capacidad vial (Daskin, 2013).

\paragraph{Demanda esperada}\label{demanda-esperada}

Para cada \(j\in\mathcal{J}\), sea \[
d_j\in\mathbb{R}_{++}
\] la \textbf{demanda esperada} en la zona \(j\) durante el horizonte
considerado. Adicionalmente definimos el vector
\(d=(d_j)_{j\in\mathcal{J}}\). En la formulación determinista de
referencia usamos \(d_j\) como estimador puntual (p.~ej. esperanza), con
la advertencia explícita de que su uso directo puede producir soluciones
frágiles en colas de distribución (Snyder, 2006). Notación:

\[
d_j>0,\quad \forall j,\qquad d\in\mathbb{R}_{++}^n.
\]

\paragraph{Capacidades de
almacenamiento}\label{capacidades-de-almacenamiento}

Para cada \(i\in\mathcal{I}\) definimos la \textbf{capacidad máxima}
\(s_i^{\max}\) mediante \[
s_i^{\max}\in\mathbb{R}_{+}\cup\{+\infty\}.
\]

La restricción de capacidad se escribe, asociada a la lógica de
activación \(x_i\):

\[
\sum_{j\in\mathcal{J}} y_{ij} \le s_i^{\max} x_i,\qquad \forall i\in\mathcal{I}. \tag{2.5}
\]

En ausencia de un tope operativo se toma \(s_i^{\max}=+\infty\) (modelo
no capacitado); en práctica real \(s_i^{\max}\) suele venir de
disponibilidad física, limitaciones regulatorias o logísticas.

\begin{center}\rule{0.5\linewidth}{0.5pt}\end{center}

\paragraph{Propiedades y observaciones
formales}\label{propiedades-y-observaciones-formales}

\begin{enumerate}
\def\labelenumi{\arabic{enumi}.}
\item
  \textbf{Dominios y no ambigüedad:} Todos los parámetros deben venir
  acompañados de su dominio y unidades. Esto evita ambigüedades en
  pruebas de existencia, estabilidad numérica y escalado de variables
  (Lewis, 2003).
\item
  \textbf{Sensibilidad estructural:} Pequeñas variaciones en \(f_i\),
  \(c_{ij}\), \(d_j\) o \(s_i^{\max}\) pueden cambiar la estructura del
  óptimo (p.~ej. la configuración \(x^*\)). Por ello, el análisis de
  sensibilidad y escenarios es parte inseparable de la modelación
  robusta (Snyder, 2006; Shapiro et al., 2021).
\item
  \textbf{Elección de \(M\) en big-M:} Si se usa la formulación
  \(y_{ij}\le M x_i\) para acoplar variables (versus la formulación con
  \(s_i^{\max}\)), se recomienda elegir \(M\) igual a una cota realista
  (ej. \(\sum_j d_j\)) para evitar degradación numérica; empero, la
  formulación (2.5) es preferible por su interpretación física y por
  estrechar el dominio factible.
\end{enumerate}

\begin{center}\rule{0.5\linewidth}{0.5pt}\end{center}

\paragraph{Ejemplo numérico paso a paso (verificación de factibilidad y
cálculo de
costo)}\label{ejemplo-numuxe9rico-paso-a-paso-verificaciuxf3n-de-factibilidad-y-cuxe1lculo-de-costo}

Tomemos la instancia didáctica ya utilizada en las secciones previas:

\begin{itemize}
\tightlist
\item
  \(\mathcal{I}=\{1,2,3\}\), \(\mathcal{J}=\{1,2,3,4\}\).
\item
  Costos fijos (miles USD): \(f=(300,400,250)\).
\item
  Matriz \(C=[c_{ij}]\) (USD/unidad):
\end{itemize}

\[
C=
\begin{bmatrix}
8 & 20 & 15 & 25 \\
12 & 18 & 10 & 30 \\
20 & 10 & 12 & 18
\end{bmatrix}.
\]

\begin{itemize}
\tightlist
\item
  Demanda: \(d=(40,20,40,25)\) (unidades).
\item
  Capacidades: \(s^{\max}=(100,80,90)\).
\item
  Consideramos la decisión \(x=(1,0,1)\) y la asignación \[
  \begin{aligned}
  &y_{11}=30,\; y_{12}=20,\; y_{13}=0,\; y_{14}=0,\\
  &y_{21}=y_{22}=y_{23}=y_{24}=0,\\
  &y_{31}=10,\; y_{32}=0,\; y_{33}=40,\; y_{34}=25.
  \end{aligned}
  \]
\end{itemize}

\textbf{Paso 1 --- Verificación de capacidades (ecuación (2.5)):}

\begin{itemize}
\tightlist
\item
  Para \(i=1\): \[
  \sum_j y_{1j}=30+20+0+0=50 \le s_1^{\max} x_1 = 100\cdot 1 = 100\quad\checkmark.
  \]
\item
  Para \(i=2\): \[
  \sum_j y_{2j}=0 \le 80\cdot 0 = 0\quad\checkmark.
  \]
\item
  Para \(i=3\): \[
  \sum_j y_{3j}=10+0+40+25=75 \le 90\cdot 1 = 90\quad\checkmark.
  \]
\end{itemize}

\textbf{Paso 2 --- Verificación de cobertura de demanda:}

\begin{itemize}
\tightlist
\item
  \(j=1\): \(y_{11}+y_{31}=30+10=40 = d_1\),
\item
  \(j=2\): \(y_{12}=20 = d_2\),
\item
  \(j=3\): \(y_{33}=40 = d_3\),
\item
  \(j=4\): \(y_{34}=25 = d_4\).
\end{itemize}

Cobertura completa \(\Rightarrow\) factible.

\textbf{Paso 3 --- Cálculo del costo total}

La función de costo (determinista, sin mantenimiento de inventario
explícito) se toma como

\[
Z(x,y)=\sum_{i\in\mathcal{I}} f_i x_i + \sum_{i\in\mathcal{I}}\sum_{j\in\mathcal{J}} c_{ij} y_{ij}.
\]

Sustituyendo:

\begin{itemize}
\tightlist
\item
  Término fijo: \(f_1 x_1 + f_3 x_3 = 300 + 250 = 550\) (miles USD).
\item
  Transporte (USD): calcular cada contribución \[
  \begin{aligned}
  &8\cdot 30 = 240,\quad 20\cdot 20 = 400,\quad 20\cdot 10 = 200,\\
  &12\cdot 40 = 480,\quad 18\cdot 25 = 450.
  \end{aligned}
  \] Suma transporte = \(240+400+200+480+450 = 1770\) USD.
\end{itemize}

Total (homogeneizando unidades a miles USD): \[
Z = 550\ (\text{miles}) + 1.770\ (\text{miles}) = 2.320\ \text{miles USD}.
\]

\textbf{Interpretación:} la solución es factible y el costo
cuantificado; variaciones en \(d\), \(c_{ij}\) o \(f_i\) deben
re-evaluarse mediante análisis de sensibilidad y, en contextos reales,
mediante formulaciones estocásticas o robustas (Snyder, 2006; Shapiro et
al., 2021).

La especificación explícita de \(f_i\), \(c_{ij}\), \(d_j\) y
\(s_i^{\max}\) permite no solo la construcción del problema determinista
base sino también su extensión a modelos estocásticos (reemplazando
\(d_j\) por variables aleatorias \(\tilde d_j\)) o robustos (definiendo
conjuntos de incertidumbre para \(d_j\) y \(c_{ij}\)). La práctica
recomendada (Lewis, 2003; Shapiro et al., 2021) es: (i) documentar
fuentes y unidades de cada parámetro; (ii) fijar cotas realistas para
\(s_i^{\max}\) o \(M\); y (iii) efectuar pruebas de sensibilidad
sistemáticas antes de adoptar políticamente cualquier solución.




\end{document}
